\chapter{Serge Petrovich Novikov --- Fields Medal 1970}\label{ch1970sergenovikov:serge-petrovich-novikov-fields-medal-1970-a-comprehensive-report-on-his-work-and-impact}

\begin{center}\rule{0.5\linewidth}{0.5pt}\end{center}

\textbf{Prepared:} July 2026 \textbf{Subject:} Serge Petrovich Novikov
(b. March 20, 1938, Gorky, USSR; d.~June 6, 2024) \textbf{Position:}
Professor, Moscow State University; University of Maryland
\textbf{Doctoral Advisor:} Mikhail Postnikov (Steklov Mathematical
Institute) \textbf{Fields Medal Citation:} ``For his contributions to
algebraic topology, especially his classification of simply connected
manifolds.''

\begin{center}\rule{0.5\linewidth}{0.5pt}\end{center}

\section{Table of Contents}\label{ch1970sergenovikov:table-of-contents}

\begin{enumerate}
\def\labelenumi{\arabic{enumi}.}
\item
 \hyperref[introduction]{Introduction}
\item
 \hyperref[biographical-background]{Biographical Background}
\item
 \hyperref[topology-and-mathematical-physics-overview]{Topology and
 Mathematical Physics: Overview}
\item
 \hyperref[cobordism-theory-and-characteristic-classes]{Cobordism
 Theory and Characteristic Classes}
\item
 \hyperref[the-adams-novikov-spectral-sequence]{The Adams-Novikov
 Spectral Sequence}
\item
 \hyperref[classification-of-simply-connected-manifolds]{Classification
 of Simply Connected Manifolds}
\item
 \hyperref[the-novikov-conjecture-on-homotopy-invariance-of-higher-signatures]{The
 Novikov Conjecture on Homotopy Invariance of Higher Signatures}
\item
 \hyperref[foliations-of-manifolds-and-the-novikov-theorem]{Foliations
 of Manifolds and the Novikov Theorem}
\item
 \hyperref[integrable-systems-and-soliton-equations]{Integrable Systems
 and Soliton Equations}
\item
 \hyperref[the-novikov-veselov-equations]{The Novikov-Veselov
 Equations}
\item
 \hyperref[schrodinger-operators-in-magnetic-fields-and-bloch-theory]{Schrodinger
 Operators in Magnetic Fields and Bloch Theory}
\item
 \hyperref[symplectic-geometry-and-multivalued-hamiltonians]{Symplectic
 Geometry and Multivalued Hamiltonians}
\item
 \hyperref[morse-theory-and-the-novikov-complex]{Morse Theory and the
 Novikov Complex}
\item
 \hyperref[quasiperiodic-structures-and-mathematical-physics]{Quasiperiodic
 Structures and Mathematical Physics}
\item
 \hyperref[quantum-field-theory-and-topological-methods]{Quantum Field
 Theory and Topological Methods}
\item
 \hyperref[unifying-techniques-in-algebraic-topology-and-physics]{Unifying
 Techniques in Algebraic Topology and Physics}
\item
 \hyperref[broad-impact-on-modern-mathematics]{Broad Impact on Modern
 Mathematics}
\item
 \hyperref[awards-and-professional-recognition]{Awards and Professional
 Recognition}
\item
 \hyperref[selected-major-publications]{Selected Major Publications}
\item
 \hyperref[ongoing-legacy-and-future-directions]{Ongoing Legacy and
 Future Directions}
\item
 \hyperref[conclusion]{Conclusion}
\item ~
 \section{\texorpdfstring{\hyperref[references-and-further-reading]{References
 and Further
 Reading}}{References and Further Reading}}\label{ch1970sergenovikov:references-and-further-reading}
\end{enumerate}

\section{1. Introduction}\label{ch1970sergenovikov:introduction}

Serge Petrovich Novikov stands as one of the most towering figures in
the mathematics of the second half of the twentieth century. His
scientific work is characterized by an substantial breadth that
bridged the deep abstractions of algebraic topology and the concrete
dynamics of mathematical physics. Awarded the Fields Medal in 1970 at
the age of thirty-two, Novikov was initially recognized for his
ground-breaking discoveries in the topology of manifolds. His early work
reshaped cobordism theory, established the stable homotopy theory of
spheres on a new foundation, solved the classification problem for
simply connected manifolds of high dimensions, and resolved
long-standing questions regarding foliation theory on 3-manifolds.

However, Novikov's contributions did not stop at pure geometry. In the
early 1970s, he made a conscious and dramatic transition toward
mathematical physics. He realized that the sophisticated topological
invariants he and his contemporaries had developed could be applied to
solve non-linear partial differential equations, explain the quantum
properties of solids, and formalize the geometric structures underlying
quantum field theory. Through this transition, Novikov became a pioneer
in the study of integrable systems, algebro-geometric integration of
soliton equations, periodic Schrödinger operators in magnetic fields,
and multivalued Hamiltonian mechanics.

This report offers a comprehensive and exhaustive study of Serge
Petrovich Novikov's mathematical legacy. It traces his biography,
analyzes his fundamental contributions across topology and physics, and
outlines the broad, continuing influence of his ideas. By examining both
the classical topological achievements that earned him the Fields Medal
and the mathematical physics research that occupied his later decades,
we highlight the deep conceptual unity that defined his scientific
career.

\begin{center}\rule{0.5\linewidth}{0.5pt}\end{center}

\section{2. Biographical Background}\label{ch1970sergenovikov:biographical-background}

\subsection{2.1 Family Heritage and Early
Life}\label{ch1970sergenovikov:family-heritage-and-early-life}

Serge Petrovich Novikov was born on March 20, 1938, in Gorky (now Nizhny
Novgorod), Soviet Union. He was born into what can only be described as
a mathematical dynasty. His father, Petr Sergeyevich Novikov
(1901--1975), was an outstanding mathematician and academician, widely
well-known for his work in mathematical logic, set theory, and group
theory, most notably his 1955 proof of the algorithmic unsolvability of
the word problem for finitely presented groups. His mother, Elizaveta
Alexandrovna Kozlova (1900--1996), was a distinguished mathematician in
her own right, specializing in descriptive set theory and topology, and
was one of the early students of Pavel Alexandrov.

This rich intellectual environment extended beyond his immediate
parents. His maternal uncle, Petr Kapitsa, was a world-famous
experimental physicist who discovered superfluidity and was awarded the
Nobel Prize in Physics in 1978. Growing up in a household where
scientific debates were a daily occurrence, the young Novikov developed
an early appreciation for the natural sciences. Although initially
interested in physics, his focus shifted toward mathematics during his
teenage years, guided by the influence of his parents and the prominent
scientists who regularly visited their home.

\subsection{2.2 Education at Moscow State
University}\label{ch1970sergenovikov:education-at-moscow-state-university}

In 1955, Novikov enrolled in the Mechanics and Mathematics Faculty
(Mech-Mat) of Moscow State University (MSU). This period marked the
absolute golden age of Soviet mathematics. The faculty at Mech-Mat
included notable figures such as Andrey Kolmogorov, Israel Gelfand,
Lev Pontryagin, Pavel Alexandrov, and Ivan Petrovsky. Novikov was
quickly swept up in this atmosphere of intense mathematical activity.

He graduated from MSU in 1960 and immediately began his postgraduate
studies at the Steklov Mathematical Institute under the supervision of
Mikhail Postnikov, a leading Soviet topologist. Postnikov provided
Novikov with a rigorous training in algebraic topology, particularly
homotopy theory. Novikov completed his Candidate of Sciences degree
(equivalent to a PhD) in 1964, and just one year later, in 1965, he
defended his Doctor of Sciences dissertation (equivalent to the
habilitation). His dissertation, which focused on the topology of
manifolds and cobordism theory, established him as one of the leading
topologists of his generation.

\subsection{2.3 Academic Career and Academic
Migration}\label{ch1970sergenovikov:academic-career-and-academic-migration}

Novikov began his professional career as a researcher at the Steklov
Mathematical Institute, working there from 1963 to 1971. In 1967, he was
appointed a professor at Moscow State University, a position he held
concurrently with his research post. In 1971, he moved his primary
research affiliation to the Landau Institute for Theoretical Physics of
the USSR Academy of Sciences, a transition that reflected his growing
interest in mathematical physics.

In 1983, Novikov was appointed Head of the Department of Higher Geometry
and Topology at Moscow State University, succeeding his former mentors.
He used this position to build a major research school, integrating
topology and physics. Following the collapse of the Soviet Union and the
subsequent opening of borders, Novikov began to spend significant time
in the West. In 1996, he joined the faculty of the University of
Maryland, College Park, as a distinguished university professor,
splitting his academic year between Maryland and Moscow State
University. Novikov remained active in research and teaching until his
death on June 6, 2024, at the age of eighty-six.

\begin{center}\rule{0.5\linewidth}{0.5pt}\end{center}

\section{3. Topology and Mathematical Physics:
Overview}\label{ch1970sergenovikov:topology-and-mathematical-physics-overview}

\subsection{3.1 The Great Bifurcation and Conceptual
Unity}\label{ch1970sergenovikov:the-great-bifurcation-and-conceptual-unity}

To understand Novikov's work, one must examine the apparent division of
his career into two phases: algebraic topology (1959--1971) and
mathematical physics (1971--2024). In the early phase, Novikov was a
pure topologist, dealing with Thom complexes, homotopy groups, and
surgery theory. In the later phase, he focused on Schrödinger operators,
the Korteweg-de Vries equation, and Fermi surfaces.

However, Novikov himself always emphasized the fundamental unity of
these two phases. He did not view his transition to physics as an
abandonment of topology, but rather as an extension of it. He believed
that physical systems provide the most natural geometric structures, and
that the abstract topological invariants of pure mathematics are the
keys to unlocking the behavior of physical systems. For Novikov,
topology was not just an algebraic game; it was the study of the
qualitative, global properties of space, which manifests directly in the
physical laws governing the universe.

\subsection{3.2 The Moscow School and the Novikov
Seminar}\label{ch1970sergenovikov:the-moscow-school-and-the-novikov-seminar}

One of Novikov's greatest achievements was the creation of a unique
research school that united pure topologists and theoretical physicists.
The centerpiece of this school was the Novikov Seminar, which met weekly
at Moscow State University. Unlike traditional seminars that focused on
narrow subfields, Novikov's seminar was highly interdisciplinary.

In this seminar, topologists were forced to learn the language of
quantum mechanics and hydrodynamics, while physicists were introduced to
cobordism, characteristic classes, and K-theory. This cross-pollination
led to several major discoveries. The ``Novikov school'' trained a
generation of productive mathematical physicists---such as Boris
Dubrovin, Igor Krichever, Aleksandr Veselov, Victor Buchstaber, and
others---who went on to play leading roles in the development of modern
integrable systems, topological quantum field theory, and string theory.

\begin{center}\rule{0.5\linewidth}{0.5pt}\end{center}

\section{4. Cobordism Theory and Characteristic
Classes}\label{ch1970sergenovikov:cobordism-theory-and-characteristic-classes}

\subsection{4.1 Thom's Foundation and Novikov's
Intervention}\label{ch1970sergenovikov:thoms-foundation-and-novikovs-intervention}

Cobordism theory was initiated by René Thom in his path-breaking 1954
paper. Two closed, smooth \(n\)-dimensional manifolds \(M_1\) and
\(M_2\) are said to be cobordant if there exists a smooth, compact
\((n+1)\)-dimensional manifold \(W\) with boundary
\(\partial W = M_1 \sqcup M_2\). Cobordism defines an equivalence
relation, and the set of cobordism classes of \(n\)-dimensional
manifolds forms an abelian group \(\Omega_n\). Under the operations of
disjoint union (addition) and Cartesian product (multiplication), the
direct sum of these groups forms a graded ring: \[
\Omega_* = \bigoplus_{n \ge 0} \Omega_n
\] Thom used the homotopy groups of special spaces, now called Thom
complexes, to compute the unoriented cobordism ring.

Novikov took this theory further by analyzing cobordism groups with
additional structures on the tangent or normal bundles. He focused
particularly on complex cobordism (unitary cobordism, denoted \(MU_*\)
or \(\Omega_*^U\)) and symplectic cobordism (\(MSp_*\)). Manifolds
representing classes in \(MU_*\) have a stable complex structure on
their normal bundle, which connects them directly to complex algebraic
geometry.

\subsection{4.2 Calculation of the Unitary Cobordism
Ring}\label{ch1970sergenovikov:calculation-of-the-unitary-cobordism-ring}

In a series of papers in the early 1960s, Novikov computed the ring
structure of complex cobordism. He proved that the complex cobordism
ring is a polynomial algebra over the integers: \[
MU_* \cong \mathbb{Z}[x_1, x_2, x_3, \dots]
\] where the generator \(x_i\) has dimension \(2i\). This result was
mathematically significant because it showed that complex cobordism has
no torsion, making it an effective tool for cohomology calculations.

To prove this, Novikov developed new techniques in homological algebra
and homotopy theory. He analyzed the stable homotopy category of Thom
spaces and used the Adams spectral sequence to compute the stable
homotopy groups. He also computed the symplectic cobordism ring
\(MSp_*\), which proved to be much more complex due to the presence of
torsion elements, revealing deep connections to the quaternionic
structures of manifolds.

\subsection{4.3 Cobordism as a Cohomology Theory and Characteristic
Classes}\label{ch1970sergenovikov:cobordism-as-a-cohomology-theory-and-characteristic-classes}

Novikov realized that cobordism could be used to construct a generalized
cohomology theory, \(MU^*(X)\), which satisfies all the
Eilenberg-Steenrod axioms except the dimension axiom. In this
generalized theory, the cohomology of a point is not the field of
coefficients, but the complex cobordism ring \(MU_*\).

Novikov used this framework to study characteristic classes. Classical
characteristic classes, such as Chern and Pontryagin classes, take
values in ordinary cohomology \(H^*(X; \mathbb{Z})\). Novikov introduced
generalized characteristic classes taking values in \(MU^*(X)\). He
showed that these generalized classes contain far more topological
information than classical ones, and he used them to establish new
relations between Pontryagin classes on differentiable manifolds, laying
the groundwork for the classification of smooth structures.

\begin{center}\rule{0.5\linewidth}{0.5pt}\end{center}

\section{5. The Adams-Novikov Spectral
Sequence}\label{ch1970sergenovikov:the-adams-novikov-spectral-sequence}

\subsection{5.1 The Classical Adams Spectral Sequence and its
Limitations}\label{ch1970sergenovikov:the-classical-adams-spectral-sequence-and-its-limitations}

In stable homotopy theory, computing the stable homotopy groups of
spheres, \(\pi_*^s\), is one of the most fundamental and difficult
problems. In 1958, J. Frank Adams introduced a powerful algebraic tool,
the Adams spectral sequence, which computes these groups using the
homological properties of ordinary cohomology over the mod \(p\)
Steenrod algebra: \[
E_2^{s,t} = \operatorname{Ext}_{\mathcal{A}_p}^{s,t}(\mathbb{Z}_p, \mathbb{Z}_p) \implies (\pi_{t-s}^s)_p^\wedge
\] where \(\mathcal{A}_p\) is the mod \(p\) Steenrod algebra.

While the classical Adams spectral sequence was highly successful, it
had major limitations. The algebra \(\mathcal{A}_p\) is extremely
complex, making the calculation of the \(E_2\)-term difficult beyond low
dimensions. Furthermore, there are many non-trivial differentials
(\(d_r\) for \(r \ge 2\)) and extension problems that make finding the
actual homotopy groups highly challenging.

\subsection{5.2 The Breakthrough: Complex Cobordism as the Cohomology
Theory}\label{ch1970sergenovikov:the-breakthrough-complex-cobordism-as-the-cohomology-theory}

In his landmark 1967 paper, \emph{``The methods of algebraic topology
from the viewpoint of cobordism theory,''} Novikov proposed a radical
modification of the Adams spectral sequence. He suggested replacing
ordinary cohomology \(H^*(X; \mathbb{Z}_p)\) with complex cobordism
\(MU^*(X)\).

This replacement yielded what is now called the \textbf{Adams-Novikov
Spectral Sequence} (ANSS). The spectral sequence has the form: \[
E_2^{s,t} = \operatorname{Ext}_{MU_*(MU)}^{s,t}(MU_*, MU_*) \implies \pi_{t-s}^s
\] where \(MU_*(MU)\) is the algebra of cooperations of the complex
cobordism theory. Because \(MU\) contains a vast amount of geometric and
arithmetic information, the \(E_2\)-term of the ANSS is much closer to
the actual stable homotopy groups than the \(E_2\)-term of the classical
Adams spectral sequence. In particular, the ANSS has far fewer
non-trivial differentials, and many differentials that are difficult to
compute in the classical sequence vanish automatically in the ANSS.

\begin{verbatim}
 Stable Homotopy Groups of Spheres \pi_*^s
 ^
 | (Convergence)
 +-------------------------------+
 | Adams-Novikov Spectral Seq. |
 | (ANSS) |
 +-------------------------------+
 ^
 |
 Ext_{MU_*(MU)}^{s,t}(MU_*, MU_*) <--- (E_2 Term)
 ^
 | (Structure determined by)
 +-------------------------------+
 | Novikov Operations & |
 | Formal Group Laws (Quillen) |
 +-------------------------------+
\end{verbatim}

\subsection{5.3 Novikov Operations and the Connection to Formal Group
Laws}\label{ch1970sergenovikov:novikov-operations-and-the-connection-to-formal-group-laws}

To compute the \(E_2\)-term of the ANSS, Novikov had to determine the
structure of the algebra of stable operations on complex cobordism,
\(MU^*(MU)\). He constructed a family of operations, now known as the
\textbf{Novikov operations} (or Landweber-Novikov operations), which act
on \(MU^*(X)\) and generalize the Steenrod operations.

Soon after Novikov's work, Daniel Quillen made a major discovery: he
proved that the complex cobordism ring \(MU_*\) is isomorphic to the
universal ring of formal group laws in arithmetic geometry. Under this
isomorphism, the cooperation algebra \(MU_*(MU)\) corresponds to the
ring of isomorphisms between formal group laws. This discovery linked
stable homotopy theory with algebraic number theory and algebraic
geometry, leading to the development of chromatic homotopy theory, which
classifies stable homotopy types using the height of formal groups.

\begin{center}\rule{0.5\linewidth}{0.5pt}\end{center}

\section{6. Classification of Simply Connected
Manifolds}\label{ch1970sergenovikov:classification-of-simply-connected-manifolds}

\subsection{6.1 Surgery Theory and the Reconstruction of
Manifolds}\label{ch1970sergenovikov:surgery-theory-and-the-reconstruction-of-manifolds}

In the early 1960s, differential topology was revolutionized by the
development of surgery theory, pioneered by William Browder, Sergei
Novikov, and C.T.C. Wall. Surgery is a geometric technique for modifying
a manifold to change its homotopy groups while preserving its dimension.
Given an \(n\)-dimensional manifold \(M\) and an embedding of a sphere
\(S^k \times D^{n-k} \subset M\), one cuts out the interior of this
product and glues in \(D^{k+1} \times S^{n-k-1}\) along their common
boundary \(S^k \times S^{n-k-1}\).

Surgery theory asks: given a Poincaré complex \(X\) (a topological space
satisfying Poincaré duality, which behaves homotopically like a
manifold), does there exist a smooth manifold \(M\) homotopy equivalent
to \(X\)? If so, how many such smooth manifolds exist up to
diffeomorphism?

\subsection{6.2 Novikov's Classification
Theorem}\label{ch1970sergenovikov:novikovs-classification-theorem}

Novikov used surgery theory to classify simply connected, smooth
manifolds of dimension \(n \ge 5\). He proved that for simply connected
manifolds, the obstacle to finding a smooth manifold in a given homotopy
type lies in the stable normal bundle and a set of algebraic invariants
known as surgery obstructions (which lie in the L-groups
\(L_n(\mathbb{Z})\)).

Novikov's classification theorem showed that if two simply connected,
smooth manifolds \(M_1\) and \(M_2\) of dimension \(n \ge 5\) are
homotopy equivalent, and if their rational Pontryagin classes match,
then they are diffeomorphic up to a finite number of choices. This
theorem provided a systematic way to classify manifolds, reducing a
difficult geometric problem to a set of calculations in homotopy theory
and characteristic classes.

\subsection{6.3 Homotopy Types vs.~Diffeomorphism
Types}\label{ch1970sergenovikov:homotopy-types-vs.-diffeomorphism-types}

A key consequence of Novikov's work was the clear mathematical
distinction between the homotopy category and the smooth category of
manifolds. Before this work, it was unclear how many distinct smooth
structures could exist on a single homotopy type. Novikov's results,
combined with Milnor's discovery of exotic spheres, showed that while
the homotopy type determines the coarse topology, the smooth structures
are governed by Pontryagin classes and the homological properties of the
tangent bundle. This work provided the theoretical foundation for
high-dimensional geometric topology.

\begin{center}\rule{0.5\linewidth}{0.5pt}\end{center}

\section{7. The Novikov Conjecture on Homotopy Invariance of Higher
Signatures}\label{ch1970sergenovikov:the-novikov-conjecture-on-homotopy-invariance-of-higher-signatures}

\subsection{7.1 Hirzebruch's Signature Theorem and its
Generalization}\label{ch1970sergenovikov:hirzebruchs-signature-theorem-and-its-generalization}

For a closed, oriented, smooth manifold \(M\) of dimension \(4k\), the
signature \(\sigma(M)\) is defined as the signature of the quadratic
intersection form on the middle-dimensional cohomology group
\(H^{2k}(M; \mathbb{R})\). Hirzebruch's Signature Theorem states that
the signature is a Pontryagin number: \[
\sigma(M) = \langle L_k(p_1, \dots, p_k), [M] \rangle
\] where \(L_k\) is the \(L\)-class, a polynomial in the Pontryagin
classes \(p_i\), and \([M]\) is the fundamental class of the manifold.

If \(g: M' \to M\) is an orientation-preserving homotopy equivalence,
then \(\sigma(M') = \sigma(M)\). This implies that the specific
combination of Pontryagin classes appearing in the \(L\)-class is
homotopically invariant when integrated over the manifold. However, the
individual Pontryagin classes themselves are \emph{not} homotopy
invariants, as was shown by Milnor's construction of exotic structures.

\subsection{7.2 Formulation of the Novikov
Conjecture}\label{ch1970sergenovikov:formulation-of-the-novikov-conjecture}

Novikov turned his attention to non-simply connected manifolds. Let
\(M\) be a closed, oriented manifold of dimension \(n\), and let
\(G = \pi_1(M)\) be its fundamental group. Let \(B G\) be the
classifying space of \(G\), and let \(f: M \to B G\) be the classifying
map corresponding to the universal cover of \(M\).

For any cohomology class \(x \in H^*(B G; \mathbb{Q})\), we can define
the \textbf{higher signature} of \(M\) associated with \(x\) as the
rational number: \[
\sigma_x(M) = \langle L(M) \cup f^*(x), [M] \rangle
\] where \(L(M)\) is the total Hirzebruch \(L\)-class of the tangent
bundle of \(M\).

In 1970, Novikov formulated his famous conjecture: \textgreater{}
\textbf{The Novikov Conjecture on Higher Signatures:} If \(g: M' \to M\)
is an orientation-preserving homotopy equivalence between closed,
oriented manifolds, and if \(f': M' \to B G\) is the classifying map for
\(M'\), then for every class \(x \in H^*(B G; \mathbb{Q})\), the higher
signatures are equal: \textgreater{} \[
> \langle L(M') \cup (f')^*(x), [M'] \rangle = \langle L(M) \cup f^*(x), [M] \rangle
> \]

\subsection{7.3 Interdisciplinary Impact of the
Conjecture}\label{ch1970sergenovikov:interdisciplinary-impact-of-the-conjecture}

The Novikov Conjecture has had a major impact on geometric topology,
index theory, and operator algebras. It serves as a bridge connecting: *
\textbf{Surgery Theory:} The homotopy invariance of higher signatures is
equivalent to the rational injectivity of the assembly map in L-theory:
\[
 A: H_*(B G; \mathbf{L}) \to L_*(G)
 \] * \textbf{Index Theory:} Via the Atiyah-Singer index theorem,
higher signatures can be interpreted as indices of elliptic operators
(specifically the signature operator) twisted by flat bundles associated
with the fundamental group. * **C*-Algebras and K-Theory:** The
Baum-Connes conjecture, which predicts the K-theory of the reduced group
C*-algebra \(C^*_r(G)\), implies the Novikov Conjecture.

The conjecture has been proven for many classes of groups, including: 1.
Free groups and free abelian groups (Novikov's original results). 2.
Word-hyperbolic groups (Gromov, Connes, Moscovici). 3. Discrete
subgroups of Lie groups (Kasparov). 4. Groups acting on buildings and
coarse metric spaces (Yu).

Despite this progress, the general Novikov Conjecture remains open and
continues to drive research in geometric topology and non-commutative
geometry.

\begin{center}\rule{0.5\linewidth}{0.5pt}\end{center}

\section{8. Foliations of Manifolds and the Novikov
Theorem}\label{ch1970sergenovikov:foliations-of-manifolds-and-the-novikov-theorem}

\subsection{8.1 Codimension-One
Foliations}\label{ch1970sergenovikov:codimension-one-foliations}

A foliation of codimension one on a 3-dimensional manifold \(M^3\) is a
partition of the manifold into a family of disjoint, connected,
2-dimensional submanifolds (called leaves) that locally look like
parallel planes in \(\mathbb{R}^3\). Foliations represent the geometric
realization of integrable systems of differential forms, as governed by
Frobenius's theorem.

In the 1950s, Georges Reeb constructed a smooth codimension-one
foliation of the 3-sphere \(S^3\). Reeb's foliation contained a single
compact leaf, which was a torus \(T^2\). This torus divided \(S^3\) into
two solid tori, each foliated by leaves that wound around the boundary
torus. Topologists wondered whether \emph{every} codimension-one
foliation of \(S^3\) must contain a compact leaf.

\subsection{8.2 The Novikov Foliation
Theorem}\label{ch1970sergenovikov:the-novikov-foliation-theorem}

In 1965, Novikov answered this question in the affirmative. He proved
the following theorem: \textgreater{} \textbf{Novikov's Foliation
Theorem:} Let \(M^3\) be a closed, oriented 3-manifold, and let
\(\mathcal{F}\) be a smooth (\(C^2\)) codimension-one foliation on
\(M^3\). If the fundamental group \(\pi_1(M^3)\) is finite, or if the
second homotopy group \(\pi_2(M^3) \neq 0\), then the foliation
\(\mathcal{F}\) contains at least one compact leaf which is a torus
\(T^2\) (or a sphere \(S^2\) in some degenerate cases). In particular,
any codimension-one foliation of \(S^3\) contains a compact torus leaf.

\subsection{8.3 Proof Strategy and
Legacy}\label{ch1970sergenovikov:proof-strategy-and-legacy}

Novikov's proof introduced several key concepts to foliation theory: *
\textbf{Transverse Loops:} He studied closed loops that are everywhere
transverse to the leaves of the foliation. * \textbf{Vanishing Cycles:}
He showed that if a foliation lacks compact leaves, one can construct a
sequence of loops on the leaves that contract in the ambient manifold
but not within the leaves themselves, leading to a contradiction. *
\textbf{Reeb Components:} He proved that the compact torus leaf must
bound a solid torus foliated in the manner of the Reeb foliation, now
called a Reeb component or Novikov component.

Novikov's theorem was a milestone in the qualitative theory of
foliations. It showed that the global topological structure of a
manifold can force the existence of localized, compact geometric
features in any foliation, a principle that remains central to modern
3-manifold topology.

\begin{center}\rule{0.5\linewidth}{0.5pt}\end{center}

\section{9. Integrable Systems and Soliton
Equations}\label{ch1970sergenovikov:integrable-systems-and-soliton-equations}

\subsection{9.1 The Soliton Revolution and the Periodic KdV
Equation}\label{ch1970sergenovikov:the-soliton-revolution-and-the-periodic-kdv-equation}

In 1967, Gardner, Greene, Kruskal, and Miura (GGKM) discovered the
Inverse Scattering Transform (IST) to solve the Korteweg-de Vries (KdV)
equation: \[
u_t - 6u u_x + u_{xxx} = 0
\] for initial data \(u(x, 0)\) that decay rapidly as
\(x \to \pm \infty\). The IST is a non-linear analog of the Fourier
transform: it associates the solution \(u(x, t)\) with the potential of
a one-dimensional Schrödinger operator: \[
L(t) = -\frac{d^2}{dx^2} + u(x, t)
\] As \(u(x, t)\) evolves according to the KdV equation, the spectrum of
\(L(t)\) remains constant.

However, the GGKM method relied on the scattering data of \(L(t)\) on
the entire real line, which is not defined for periodic initial data
\(u(x + T, t) = u(x, t)\). Solving the periodic KdV equation was a major
open problem in the early 1970s.

\subsection{9.2 Algebro-Geometric (Finite-Gap)
Integration}\label{ch1970sergenovikov:algebro-geometric-finite-gap-integration}

Novikov solved the periodic KdV problem in 1974 by introducing the
method of \textbf{finite-gap integration}. He looked for potentials
\(u(x)\) for which the Schrödinger operator
\(L = -\frac{d^2}{dx^2} + u(x)\) has a finite number of gaps
(instability intervals) in its spectrum. For a generic periodic
potential, the spectrum consists of an infinite sequence of bands
separated by gaps. Novikov studied the special case where all but a
finite number of gaps vanish.

Novikov proved that these finite-gap potentials satisfy a system of
ordinary differential equations (now called the Novikov equations),
which represent stationary higher KdV equations. He showed that the set
of finite-gap potentials with a fixed spectrum forms a Jacobi variety of
a hyperelliptic Riemann surface (the spectral curve). The evolution of
the potential under the KdV flow corresponds to linear motion on this
Jacobian torus.

\begin{verbatim}
 Periodic KdV Solution u(x,t)
 ^
 | (Evaluated via Jacobi theta functions)
 +---------------------------------------------+
 | Abel-Jacobi Map on the Jacobian |
 | J(C) of the Spectral Curve C |
 +---------------------------------------------+
 ^
 | (Linear flow)
 Flow on the Jacobian: z(t) = z(0) + V t
 ^
 | (Defined by)
 Riemann Surface y^2 = \prod_{i=1}^{2g+1} (E - E_i)
 with branch points at spectral gap edges E_i
\end{verbatim}

\subsection{9.3 Theta Function Formulae and the Dubrovin-Krichever-Its
Synthesis}\label{ch1970sergenovikov:theta-function-formulae-and-the-dubrovin-krichever-its-synthesis}

Working with his collaborators Boris Dubrovin, Igor Krichever, and
Alexander Its, Novikov expressed the solutions to the periodic KdV
equation explicitly using Riemann theta functions. The formula for the
solution is: \[
u(x, t) = -2 \frac{\partial^2}{\partial x^2} \ln \theta(U x + W t + z_0; \tau) + C
\] where \(\theta\) is the multidimensional theta function associated
with the spectral curve, \(U\) and \(W\) are vectors of periods of
meromorphic differentials, and \(\tau\) is the period matrix of the
Riemann surface.

This algebro-geometric method was soon extended to a wide variety of
non-linear wave equations, including the non-linear Schrödinger
equation, the Sine-Gordon equation, and the KP equation. Krichever
generalized this approach by introducing ``Krichever data'' (a Riemann
surface, a point on it, and a local coordinate), which became a standard
framework for studying integrable systems.

\begin{center}\rule{0.5\linewidth}{0.5pt}\end{center}

\section{10. The Novikov-Veselov
Equations}\label{ch1970sergenovikov:the-novikov-veselov-equations}

\subsection{10.1 Generalization of KdV to 2+1
Dimensions}\label{ch1970sergenovikov:generalization-of-kdv-to-21-dimensions}

Following the success of the finite-gap integration method for 1D
systems, Novikov and his student Aleksandr Veselov sought to extend
these results to higher dimensions. They focused on
\((2+1)\)-dimensional systems, where the spatial variables are complex
coordinates \(z = x + iy\) and \(\bar{z} = x - iy\).

In 1984, they introduced the \textbf{Novikov-Veselov (NV) equation} for
a real-valued function \(v(z, \bar{z}, t)\): \[
\partial_t v = 8 \operatorname{Re}(\partial^3 v) - 3 \partial(v \partial^{-1} \bar{\partial} v) - 3 \bar{\partial}(v \bar{\partial}^{-1} \partial v)
\] where the operators are defined as: \[
\partial = \frac{1}{2}\left( \frac{\partial}{\partial x} - i\frac{\partial}{\partial y} \right), \quad \bar{\partial} = \frac{1}{2}\left( \frac{\partial}{\partial x} + i\frac{\partial}{\partial y} \right)
\] and \(\partial^{-1}\), \(\bar{\partial}^{-1}\) represent the inverse
operators (integration).

\subsection{10.2 Connection to the 2D Schrödinger
Operator}\label{ch1970sergenovikov:connection-to-the-2d-schruxf6dinger-operator}

The NV equation is integrable via the Lax representation, which in this
case is a Manakov-type triad: \[
\frac{\partial L}{\partial t} = [A, L] + B L
\] where \(L\) is the two-dimensional Schrödinger operator: \[
L = -\Delta + v(z, \bar{z}, t) = -4 \partial \bar{\partial} + v(z, \bar{z}, t)
\] Crucially, this Lax relation is valid only at a fixed energy level
\(E\), meaning we look for solutions to \(L \psi = E \psi\). The NV
equation governs the deformations of the potential \(v\) that preserve
the fixed-energy spectral data.

\subsection{10.3 Algebro-Geometric Solutions and Physical
Realization}\label{ch1970sergenovikov:algebro-geometric-solutions-and-physical-realization}

Veselov and Novikov constructed algebro-geometric solutions for the NV
equation. The spectral data for the 2D Schrödinger operator at a fixed
energy are represented by a Riemann surface (the ``spectral curve'')
with a set of points where the wave functions have poles.

The NV equation provides a model for studying 2D wave propagation in
dispersive media. Mathematically, it remains a key example of how the
spectral theory of multi-dimensional operators can be integrated using
Riemann surfaces, theta functions, and algebraic geometry.

\begin{center}\rule{0.5\linewidth}{0.5pt}\end{center}

\section{11. Schrodinger Operators in Magnetic Fields and Bloch
Theory}\label{ch1970sergenovikov:schrodinger-operators-in-magnetic-fields-and-bloch-theory}

\subsection{11.1 The Quantum Mechanics of Periodic Magnetic
Fields}\label{ch1970sergenovikov:the-quantum-mechanics-of-periodic-magnetic-fields}

In solid-state physics, Bloch's theorem describes the quantum states of
an electron in a periodic crystal potential \(V(\mathbf{x})\). The wave
functions, called Bloch waves, satisfy: \[
\psi_{\mathbf{k}}(\mathbf{x} + \mathbf{a}) = e^{i \mathbf{k} \cdot \mathbf{a}} \psi_{\mathbf{k}}(\mathbf{x})
\] where \(\mathbf{a}\) is a lattice vector and \(\mathbf{k}\) is the
crystal momentum (quasimomentum).

In the late 1970s, Novikov studied the case where the electron is also
subject to a periodic magnetic field
\(\mathbf{B}(\mathbf{x}) = \nabla \times \mathbf{A}(\mathbf{x})\). The
Hamiltonian is: \[
H = \frac{1}{2m} \left( -i\hbar \nabla - \frac{e}{c} \mathbf{A}(\mathbf{x}) \right)^2 + V(\mathbf{x})
\] If the magnetic flux through a unit cell of the lattice is a rational
multiple of the flux quantum \(\Phi_0 = hc/e\), the magnetic translation
operators commute, and one can define a magnetic Brillouin zone.
However, if the magnetic flux is irrational, the translation operators
do not commute, and the spectrum of the operator becomes highly complex
(often forming a Cantor set, such as the Hofstadter butterfly).

\subsection{11.2 Topological Bloch Theory and Chern
Numbers}\label{ch1970sergenovikov:topological-bloch-theory-and-chern-numbers}

Novikov developed a topological formulation of Bloch theory in the
rational flux case. He showed that the Bloch wave functions
\(\psi_{\mathbf{k}}(\mathbf{x})\) define a complex vector bundle over
the Brillouin zone (which is topologically a torus \(T^2\)).

\begin{verbatim}
 Brillouin Zone Torus T^2 = S^1 x S^1
 |
 v (Associate to each point k \in T^2)
 Eigenspace of H at quasimomentum k
 |
 v
 Complex Vector Bundle E over T^2
 |
 v
 First Chern Class: c_1(E) \in H^2(T^2; Z) \cong Z
 |
 v (Physical Interpretation)
 Quantized Hall Conductivity: \sigma_{xy} = \frac{e^2}{h} c_1(E)
\end{verbatim}

Novikov proved that the first Chern class of this vector bundle,
\(c_1(E) \in H^2(T^2; \mathbb{Z}) \cong \mathbb{Z}\), is a topological
invariant. This invariant represents the Chern number of the magnetic
Bloch band. Via the Kubo formula, this Chern number is directly
proportional to the Hall conductivity: \[
\sigma_{xy} = \frac{e^2}{h} \sum_{n \text{ (occupied)}} C_n
\] where \(C_n\) is the Chern number of the \(n\)-th band. This
topological formulation provided the mathematical foundation for the
quantum Hall effect (TKNN formula) and the modern theory of topological
insulators.

\begin{center}\rule{0.5\linewidth}{0.5pt}\end{center}

\section{12. Symplectic Geometry and Multivalued
Hamiltonians}\label{ch1970sergenovikov:symplectic-geometry-and-multivalued-hamiltonians}

\subsection{12.1 Non-Exact Symplectic Forms in
Physics}\label{ch1970sergenovikov:non-exact-symplectic-forms-in-physics}

In classical mechanics, the phase space of a system is a symplectic
manifold \((M, \omega)\), where \(\omega\) is a closed, non-degenerate
2-form. The equations of motion are given by a Hamiltonian vector field
\(X_H\), defined by: \[
\iota_{X_H} \omega = dH
\] where \(H: M \to \mathbb{R}\) is the Hamiltonian function. This
formulation assumes that the class of the 1-form \(\iota_{X_H} \omega\)
is exact.

However, many physical systems do not fit this classical model. For
example, a charged particle moving in the field of a magnetic monopole
cannot be described by a globally defined vector potential
\(\mathbf{A}\). The symplectic form on the phase space is: \[
\omega = \omega_0 + e \pi^*(F)
\] where \(\omega_0\) is the standard symplectic form on the cotangent
bundle \(T^*N\), \(\pi: T^*N \to N\) is the projection, and \(F\) is the
magnetic field 2-form. If the magnetic field has a non-vanishing
monopole charge, \(F\) is closed but not exact. Thus, \(\omega\) is
closed but not exact.

\subsection{12.2 Multivalued Actions and
Hamiltonians}\label{ch1970sergenovikov:multivalued-actions-and-hamiltonians}

Novikov developed the mathematical theory of \textbf{multivalued
Hamiltonians} and Lagrangian systems with multivalued actions. He showed
that for such systems, the action functional \(S\) cannot be defined as
a single-valued function on the space of paths. Instead, the variation
of the action \(\delta S\) is a well-defined, closed 1-form on the path
space.

Novikov formulated a variational framework for these systems. He showed
that the classical trajectories correspond to the zeros of this closed
1-form. The topological invariants of the phase space, specifically the
cohomology class \([\omega] \in H^2(M; \mathbb{R})\), determine the
global properties of the trajectories. This theory provided a rigorous
geometric framework for studying particles in strong magnetic fields and
gauge fields, and it laid the foundation for the mathematical study of
anomalies in quantum field theory.

\begin{center}\rule{0.5\linewidth}{0.5pt}\end{center}

\section{13. Morse Theory and the Novikov
Complex}\label{ch1970sergenovikov:morse-theory-and-the-novikov-complex}

\subsection{13.1 Generalization of Classical Morse
Theory}\label{ch1970sergenovikov:generalization-of-classical-morse-theory}

Classical Morse theory, developed by Marston Morse, relates the topology
of a closed manifold \(M\) to the critical points of a smooth function
\(f: M \to \mathbb{R}\). If \(f\) is a Morse function (all its critical
points are non-degenerate), the Morse inequalities state that: \[
M_p(f) \ge b_p(M)
\] where \(M_p(f)\) is the number of critical points of index \(p\), and
\(b_p(M)\) is the \(p\)-th Betti number of \(M\). The homology of \(M\)
can be computed using the Morse complex: \[
\cdots \to C_{p+1} \xrightarrow{\partial} C_p \xrightarrow{\partial} C_{p-1} \to \cdots
\] where \(C_p\) is the free abelian group generated by the critical
points of index \(p\), and the boundary operator \(\partial\) counts the
gradient flow lines connecting the critical points.

In the early 1980s, Novikov generalized this theory to closed 1-forms.
Let \(\omega\) be a closed 1-form on \(M\). If \(\omega\) is not exact
(\(\omega \neq df\)), it cannot be integrated to a single-valued
function. However, locally we can write \(\omega = df_i\). The critical
points of \(\omega\) are the points where \(\omega = 0\).

\subsection{13.2 The Novikov Ring and Novikov
Homology}\label{ch1970sergenovikov:the-novikov-ring-and-novikov-homology}

To construct a homology theory associated with a closed 1-form
\(\omega\), Novikov introduced a new algebraic structure, now known as
the \textbf{Novikov ring}.

Let \(G = \pi_1(M)\) be the fundamental group of \(M\), and let
\(p: \tilde{M} \to M\) be the universal cover. The pullback
\(p^*\omega\) is exact, so \(p^*\omega = d\tilde{f}\) for a
single-valued function \(\tilde{f}: \tilde{M} \to \mathbb{R}\). The
function \(\tilde{f}\) is not periodic, but it satisfies: \[
\tilde{f}(g \cdot x) = \tilde{f}(x) + c(g)
\] for \(g \in G\), where \(c: G \to \mathbb{R}\) is the period
homomorphism defined by the cohomology class of \(\omega\).

The \textbf{Novikov ring} \(\Lambda\) associated with the homomorphism
\(c: G \to \mathbb{R}\) is the set of formal power series of the form:
\[
\lambda = \sum_{g \in G} a_g g
\] with coefficients \(a_g \in \mathbb{Z}\), satisfying the following
finiteness condition: for any real number \(C\), the set of elements
\(g \in G\) such that \(a_g \neq 0\) and \(c(g) > C\) is finite.

Novikov defined the \textbf{Novikov complex} as a chain complex of free
modules over the Novikov ring \(\Lambda\): \[
C_p(\omega) = \bigoplus_{\omega(x)=0, \operatorname{ind}(x)=p} \Lambda \cdot \tilde{x}
\] where \(\tilde{x}\) is a lift of the critical point \(x\) to the
universal cover. The boundary operator \(\partial\) is defined by
counting the gradient flow lines of \(\tilde{f}\) on \(\tilde{M}\).
Novikov proved that the homology of this complex, known as the
\textbf{Novikov homology} \(H_*(M; \omega)\), depends only on the
cohomology class of \(\omega\).

\subsection{13.3 Morse-Novikov
Inequalities}\label{ch1970sergenovikov:morse-novikov-inequalities}

Novikov proved the generalized \textbf{Morse-Novikov inequalities}: \[
M_p(\omega) \ge b_p(M; \omega)
\] where \(M_p(\omega)\) is the number of critical points of the closed
1-form \(\omega\) of index \(p\), and \(b_p(M; \omega)\) is the rank of
the Novikov homology module \(H_p(M; \omega)\) over the Novikov ring
\(\Lambda\).

Morse-Novikov theory has become a key tool in modern symplectic geometry
and topology. It serves as the model for Floer homology, where the
action functional is a closed, non-exact 1-form on an
infinite-dimensional loop space. The Novikov ring is essential for
defining the Floer homology of Lagrangian submanifolds and the quantum
cohomology of symplectic manifolds (Fukaya categories).

\begin{center}\rule{0.5\linewidth}{0.5pt}\end{center}

\section{14. Quasiperiodic Structures and Mathematical
Physics}\label{ch1970sergenovikov:quasiperiodic-structures-and-mathematical-physics}

\subsection{14.1 Semiclassical S.P. Novikov's Problem on Fermi
Surfaces}\label{ch1970sergenovikov:semiclassical-s.p.-novikovs-problem-on-fermi-surfaces}

In 1982, Novikov formulated a geometric problem motivated by the physics
of metals. In the semiclassical description of an electron in a 3D
crystal under a constant magnetic field \(\mathbf{H}\), the electron's
motion in momentum space (k-space) is governed by the equations: \[
\hbar \frac{d\mathbf{k}}{dt} = -\frac{e}{c} \mathbf{v}(\mathbf{k}) \times \mathbf{H}, \quad \mathbf{v}(\mathbf{k}) = \frac{1}{\hbar} \nabla_{\mathbf{k}} \epsilon(\mathbf{k})
\] where \(\epsilon(\mathbf{k})\) is the dispersion relation (energy) of
the electron.

These equations imply that the electron's energy
\(\epsilon(\mathbf{k})\) and the component of its momentum parallel to
the magnetic field, \(\mathbf{k} \cdot \mathbf{H}\), are conserved.
Consequently, the trajectories of the electron in k-space are the
intersections of the Fermi surface: \[
\mathcal{F} = \{ \mathbf{k} \in \mathbb{R}^3 \mid \epsilon(\mathbf{k}) = E_F \}
\] with planes perpendicular to the magnetic field vector
\(\mathbf{H}\). The Fermi surface \(\mathcal{F}\) is periodic under the
reciprocal lattice translations, so it can be viewed as a closed surface
embedded in the 3-torus \(T^3 = \mathbb{R}^3/\mathbb{Z}^3\).

\begin{verbatim}
 Fermi Surface \mathcal{F} in \mathbb{R}^3 (Periodic)
 |
 v (Intersected with)
 Planes perpendicular to Magnetic Field H
 |
 +------------------+------------------+
 | |
 v v
 Closed Orbits (Local) Open Orbits (Global)
 | |
 v (Lead to) v (Lead to)
 Normal Magnetoresistance Anisotropic Magnetoresistance
 |
 v
 Novikov's Conjecture:
 Orbits run along straight lines
 in quasiperiodic directions
\end{verbatim}

\subsection{14.2 Classification of Trajectories and Topological
Magnetoresistance}\label{ch1970sergenovikov:classification-of-trajectories-and-topological-magnetoresistance}

Novikov posed the problem of classifying the asymptotic behavior of
these intersection trajectories. If the intersection forms closed
curves, the electron's motion is localized. However, the intersection of
a periodic surface with a plane can also produce open, infinite
trajectories.

Novikov conjectured that for generic Fermi surfaces and generic
directions of the magnetic field \(\mathbf{H}\), the open trajectories
are either: 1. Closed or periodic. 2. Open, running along a straight
line in a quasiperiodic direction.

This conjecture, now known as \textbf{Novikov's problem on semiclassical
transport}, has been proved in many cases by Ivan Dynnikov and other
members of the Novikov school. The topological properties of these open
trajectories have direct, measurable physical consequences: when open
trajectories are present, the magnetoresistance of the metal becomes
highly anisotropic, allowing physicists to map the topology of the Fermi
surface experimentally.

\begin{center}\rule{0.5\linewidth}{0.5pt}\end{center}

\section{15. Quantum Field Theory and Topological
Methods}\label{ch1970sergenovikov:quantum-field-theory-and-topological-methods}

\subsection{15.1 The Interface of Topology and Gauge
Theories}\label{ch1970sergenovikov:the-interface-of-topology-and-gauge-theories}

In the late 1970s, the discovery of instanton and monopole solutions in
Yang-Mills gauge theories revealed that these classical field
configurations are governed by algebraic topology. Novikov was one of
the first mathematicians to recognize the importance of these
developments.

He analyzed the Yang-Mills instanton solutions on \(\mathbb{R}^4\) (or
\(S^4\)), showing that the instanton number (action) is the first
Pontryagin class (or second Chern class) of the gauge bundle: \[
Q = \frac{1}{8\pi^2} \int_{S^4} \operatorname{Tr}(F \wedge F) \in \mathbb{Z}
\] where \(F\) is the gauge field strength 2-form. This topological
invariant classifies the bundle up to isomorphism, and the self-duality
equations \(F = *F\) represent the absolute minima of the action within
a given topological class.

\subsection{15.2 Topological Anomalies}\label{ch1970sergenovikov:topological-anomalies}

Novikov also studied the topological origin of anomalies in quantum
field theory. Anomalies occur when a symmetry of the classical action is
violated by the quantum measure. Novikov showed that these anomalies can
be computed using index theory and the cohomology of the gauge group,
providing a geometric explanation for the cancellation of anomalies in
physical models.

\begin{center}\rule{0.5\linewidth}{0.5pt}\end{center}

\section{16. Unifying Techniques in Algebraic Topology and
Physics}\label{ch1970sergenovikov:unifying-techniques-in-algebraic-topology-and-physics}

\subsection{16.1 The Wess-Zumino-Novikov-Witten
Model}\label{ch1970sergenovikov:the-wess-zumino-novikov-witten-model}

A key example of Novikov's contribution to mathematical physics is his
geometric formulation of the topological term in non-linear sigma
models, which led to the Wess-Zumino-Novikov-Witten (WZNW) model in
conformal field theory.

In a 2D sigma model, the fields are maps \(g: \Sigma \to G\) from a
Riemann surface \(\Sigma\) to a Lie group \(G\). The action contains a
standard kinetic term and a topological term. The topological term
cannot be written as the integral of a local Lagrangian density on
\(\Sigma\).

Novikov solved this by extending the map \(g\) to a 3-dimensional
manifold \(W\) whose boundary is \(\Sigma\) (\(\partial W = \Sigma\)).
The topological action is then defined as the integral of a closed,
invariant 3-form \(\omega\) on the group \(G\) over the image of \(W\):
\[
S_{\text{topo}}(g) = 2\pi k \int_W \tilde{g}^*(\omega)
\] where \(\tilde{g}: W \to G\) is the extended map, and \(k\) is a
coupling constant.

\subsection{16.2 Quantization of the Coupling
Constant}\label{ch1970sergenovikov:quantization-of-the-coupling-constant}

For the quantum theory to be well-defined, the path integral factor
\(e^{i S_{\text{topo}}(g)}\) must be independent of the choice of the
3-dimensional manifold \(W\) and the extension \(\tilde{g}\). If \(W_1\)
and \(W_2\) are two different extensions, they glue together along their
boundary \(\Sigma\) to form a closed 3-manifold
\(M = W_1 \cup_{\Sigma} (-W_2)\).

The difference between the two actions is: \[
S_{\text{topo}}(\tilde{g}_1) - S_{\text{topo}}(\tilde{g}_2) = 2\pi k \int_M \Phi^*(\omega)
\] where \(\Phi: M \to G\) is the combined map. For the path integral
factor to be well-defined, this difference must be an integer multiple
of \(2\pi\), which requires: \[
k \int_M \Phi^*(\omega) \in \mathbb{Z}
\] Novikov proved that this condition is satisfied if and only if the
cohomology class of \(\omega\) is integral, and the coupling constant
\(k\) is an integer (the level). This formulation unified cobordism
theory, differential forms, and quantum mechanics, becoming a standard
construction in string theory and conformal field theory.

\begin{center}\rule{0.5\linewidth}{0.5pt}\end{center}

\section{17. Broad Impact on Modern
Mathematics}\label{ch1970sergenovikov:broad-impact-on-modern-mathematics}

\subsection{17.1 The Novikov School of
Mathematics}\label{ch1970sergenovikov:the-novikov-school-of-mathematics}

Novikov's influence on modern mathematics was amplified by his role as a
teacher and mentor. He established a major research school at Moscow
State University and the Landau Institute, producing dozens of
world-class mathematicians.

His students and collaborators include: * \textbf{Boris Dubrovin:} Who
made major contributions to integrable systems, Frobenius manifolds, and
Gromov-Witten theory. * \textbf{Igor Krichever:} Who developed the
algebro-geometric theory of soliton equations and Krichever-Novikov
algebras. * \textbf{Aleksandr Veselov:} Who co-developed the
Novikov-Veselov equation and made key contributions to integrable maps
and mathematical physics. * \textbf{Victor Buchstaber:} A leading
topologist who worked with Novikov on cobordism theory and formal group
laws. * \textbf{Sabir Gusein-Zade:} Who made major contributions to
singularity theory and algebraic geometry.

\subsection{17.2 The Landmark Textbooks}\label{ch1970sergenovikov:the-landmark-textbooks}

Novikov co-authored a famous three-volume textbook series with Boris
Dubrovin and Anatoly Fomenko, titled \emph{``Modern Geometry: Methods
and Applications.''}

First published in Russian in the late 1970s and subsequently translated
into English and other languages, this series revolutionized the
teaching of geometry. It replaced the traditional coordinate-heavy
approach with a modern, coordinate-free viewpoint, integrating
differential geometry, topology, and mathematical physics. The series
remains a standard reference worldwide, introducing generations of
students to the geometric foundations of modern physics.

\begin{center}\rule{0.5\linewidth}{0.5pt}\end{center}

\section{18. Awards and Professional
Recognition}\label{ch1970sergenovikov:awards-and-professional-recognition}

\subsection{18.1 The Fields Medal (1970)}\label{ch1970sergenovikov:the-fields-medal-1970}

In 1970, the International Mathematical Union (IMU) awarded Novikov the
Fields Medal at the International Congress of Mathematicians (ICM) in
Nice, France. The medal was awarded ``for his contributions to algebraic
topology, especially his classification of simply connected manifolds.''

However, due to the political environment of the Cold War and his
signature on a letter defending a dissident mathematician, the Soviet
authorities refused to grant Novikov an exit visa. He was unable to
travel to Nice to receive the medal in person. His colleague Lev
Pontryagin accepted the medal on his behalf. Novikov was finally able to
deliver his plenary address at the next ICM in Vancouver in 1974.

\subsection{18.2 The Wolf Prize in Mathematics
(2005)}\label{ch1970sergenovikov:the-wolf-prize-in-mathematics-2005}

In 2005, Novikov was awarded the Wolf Prize in Mathematics, one of the
most prestigious awards in the field, shared with Gregory Margulis. The
citation recognized: \textgreater{} ``his fundamental and pioneering
contributions to algebraic and differential topology, and to
mathematical physics, notably the introduction of algebro-geometric
methods.''

\subsection{18.3 Academy Memberships and Other
Honors}\label{ch1970sergenovikov:academy-memberships-and-other-honors}

Novikov received many other awards and honors, including: *
\textbf{Lobaevski Prize (1981):} Awarded by the USSR Academy of
Sciences. * \textbf{USSR Academy of Sciences:} Elected a corresponding
member in 1966 (at the age of 28) and a full academician in 1981. *
\textbf{US National Academy of Sciences:} Elected a Foreign Associate in
1994. * \textbf{Pontifical Academy of Sciences:} Elected a member in
1996. * \textbf{London Mathematical Society:} Elected an Honorary
Member.

\begin{center}\rule{0.5\linewidth}{0.5pt}\end{center}

\section{19. Selected Major
Publications}\label{ch1970sergenovikov:selected-major-publications}

Novikov was a prolific author whose papers and books reshaped several
fields. Some of his most significant publications include:

\begin{enumerate}
\def\labelenumi{\arabic{enumi}.}
\tightlist
\item
 \textbf{Novikov, S. P. (1962). ``Homotopy properties of Thom
 complexes.'' \emph{Matematicheskii Sbornik}, 57(4), 407--442.}

 \begin{itemize}
 \tightlist
 \item
 This paper contains Novikov's early calculations of the unitary
 (\(MU\)) and symplectic (\(MSp\)) cobordism rings, establishing them
 as polynomial algebras and developing the necessary homological
 tools.
 \end{itemize}
\item
 \textbf{Novikov, S. P. (1964). ``Differentiable manifolds of a given
 homotopy type.'' \emph{Izvestiya Rossiiskoi Akademii Nauk. Seriya
 Matematicheskaya}, 28(2), 365--420.}

 \begin{itemize}
 \tightlist
 \item
 A foundational paper in surgery theory, in which Novikov classifies
 simply connected differentiable manifolds of dimension \(n \ge 5\)
 up to diffeomorphism.
 \end{itemize}
\item
 \textbf{Novikov, S. P. (1965). ``Foliations of codimension 1 on
 manifolds.'' \emph{Doklady Akademii Nauk SSSR}, 162, 1241--1244.}

 \begin{itemize}
 \tightlist
 \item
 This paper contains the proof of the Novikov Foliation Theorem,
 demonstrating that any codimension-one foliation of a closed
 3-manifold with finite fundamental group contains a compact torus
 leaf.
 \end{itemize}
\item
 \textbf{Novikov, S. P. (1967). ``The methods of algebraic topology
 from the viewpoint of cobordism theory.'' \emph{Izvestiya Akademii
 Nauk SSSR. Seriya Matematicheskaya}, 31(4), 855--951.}

 \begin{itemize}
 \tightlist
 \item
 The paper in which Novikov introduced the Adams-Novikov Spectral
 Sequence, replacing ordinary cohomology with complex cobordism and
 defining the Landweber-Novikov operations.
 \end{itemize}
\item
 \textbf{Novikov, S. P. (1970). ``On manifolds with free abelian
 fundamental group and their signatures.'' \emph{Izvestiya Akademii
 Nauk SSSR. Seriya Matematicheskaya}, 34(2), 253--274.}

 \begin{itemize}
 \tightlist
 \item
 This paper formulated the Novikov Conjecture on the homotopy
 invariance of higher signatures for non-simply connected manifolds.
 \end{itemize}
\item
 \textbf{Novikov, S. P. (1974). ``A periodic problem for the
 Korteweg-de Vries equation.'' \emph{Funktsional'nyi Analiz i ego
 Prilozheniya}, 8(3), 54--66.}

 \begin{itemize}
 \tightlist
 \item
 This paper launched the finite-gap algebro-geometric integration
 method for periodic soliton equations, solving the periodic KdV
 equation using Riemann surfaces and theta functions.
 \end{itemize}
\item
 \textbf{Dubrovin, B. A., Krichever, I. M., \& Novikov, S. P. (1976).
 ``The Hamiltonian formalism of algebro-geometric integration of
 non-linear equations.'' \emph{Doklady Akademii Nauk SSSR}, 229(1),
 15--18.}

 \begin{itemize}
 \tightlist
 \item
 A paper establishing the Hamiltonian structures of finite-gap
 equations and linking them to algebraic geometry.
 \end{itemize}
\item
 \textbf{Novikov, S. P. (1981). ``Multivalued functions and quadratic
 forms. An analogue of Morse theory.'' \emph{Doklady Akademii Nauk
 SSSR}, 260(1), 31--35.}

 \begin{itemize}
 \tightlist
 \item
 The paper that initiated Morse-Novikov theory, generalizing Morse
 theory to closed 1-forms and defining the Novikov ring and Novikov
 complex.
 \end{itemize}
\item
 \textbf{Veselov, A. P., \& Novikov, S. P. (1984). ``Finite-zone
 two-dimensional Schrödinger operators. Potential operators.''
 \emph{Doklady Akademii Nauk SSSR}, 279(1), 20--24.}

 \begin{itemize}
 \tightlist
 \item
 Here, the authors introduced the Novikov-Veselov equation and
 constructed its algebro-geometric solutions.
 \end{itemize}
\item
 \textbf{Dubrovin, B. A., Fomenko, A. T., \& Novikov, S. P. (1984).
 \emph{Modern Geometry: Methods and Applications}. Springer-Verlag.}

 \begin{itemize}
 \tightlist
 \item
 The three-volume textbook series that unified differential geometry,
 topology, and mathematical physics.
 \end{itemize}
\end{enumerate}

\begin{center}\rule{0.5\linewidth}{0.5pt}\end{center}

\section{20. Ongoing Legacy and Future
Directions}\label{ch1970sergenovikov:ongoing-legacy-and-future-directions}

Novikov's mathematical legacy remains active, with several key areas of
research continuing to expand:

\subsection{20.1 The Novikov Conjecture in Geometric Topology and
K-Theory}\label{ch1970sergenovikov:the-novikov-conjecture-in-geometric-topology-and-k-theory}

The Novikov Conjecture remains one of the central problems in
high-dimensional geometric topology. Modern research focuses on proving
the conjecture for increasingly complex classes of groups. The
development of coarse geometry by Guoliang Yu and the study of the
Baum-Connes assembly map have provided new tools, but a complete proof
for all discrete groups remains an active area of research.

\subsection{20.2 Morse-Novikov Theory and Floer
Homology}\label{ch1970sergenovikov:morse-novikov-theory-and-floer-homology}

In symplectic topology, Morse-Novikov theory is a foundational tool. The
Novikov ring is essential for defining the Floer homology of Lagrangian
submanifolds and the quantum cohomology of symplectic manifolds.
Researchers continue to study the properties of the Novikov complex in
the context of Fukaya categories, mirror symmetry, and the study of
dynamical systems on manifolds.

\subsection{20.3 Semiclassical Transport on Fermi
Surfaces}\label{ch1970sergenovikov:semiclassical-transport-on-fermi-surfaces}

In condensed matter physics, Novikov's problem on semiclassical
transport remains highly relevant. With the discovery of topological
insulators, Weyl semimetals, and other novel states of matter,
understanding the open and quasiperiodic trajectories of electrons on
Fermi surfaces in strong magnetic fields has become essential for
interpreting transport measurements in these materials.

\begin{center}\rule{0.5\linewidth}{0.5pt}\end{center}

\section{21. Conclusion}\label{ch1970sergenovikov:conclusion}

Serge Petrovich Novikov was a mathematical pioneer whose work bridged
pure abstraction and physical reality. His unique approach was
characterized by a powerful geometric intuition that allowed him to see
the global, topological structures underlying physical equations, and
the physical meaning of topological invariants.

Whether he was classifying high-dimensional manifolds, computing the
stable homotopy groups of spheres, proving the existence of compact
leaves in foliations, or integrating non-linear wave equations, Novikov
always sought the most elegant and conceptually unified explanation. His
transition from topology to mathematical physics showed his intellectual
courage and his refusal to be confined by artificial boundaries between
fields.

His death on June 6, 2024, marked the end of an era, but his ideas,
conjectures, and the school of mathematics he built will continue to
guide and inspire mathematicians and physicists for generations to come.

\begin{center}\rule{0.5\linewidth}{0.5pt}\end{center}

\section{22. References and Further
Reading}\label{ch1970sergenovikov:references-and-further-reading-1}

\begin{enumerate}
\def\labelenumi{\arabic{enumi}.}
\tightlist
\item
 Adams, J. F. (1958). ``On the structure and applications of the
 Steenrod algebra.'' \emph{Commentarii Mathematici Helvetici}, 32,
 180--214.
\item
 Browder, W. (1972). \emph{Surgery on Simply Connected Manifolds}.
 Springer-Verlag.
\item
 Connes, A. (1994). \emph{Noncommutative Geometry}. Academic Press.
\item
 Dubrovin, B. A., Fomenko, A. T., \& Novikov, S. P. (1992).
 \emph{Modern Geometry: Methods and Applications} (2nd ed.).
 Springer-Verlag.
\item
 Hirzebruch, F. (1966). \emph{Topological Methods in Algebraic
 Geometry}. Springer-Verlag.
\item
 Lück, Wolfgang. (2002). ``L2-Invariants: Theory and Applications to
 Geometry and K-Theory.'' \emph{Ergebnisse der Mathematik und ihrer
 Grenzgebiete}. Springer.
\item
 Novikov, S. P. (2004). ``Topology in the 20th Century: A View from
 Within.'' \emph{Russian Mathematical Surveys}, 59(5), 3--28.
\item
 Quillen, D. (1969). ``On the formal group laws of unoriented and
 complex cobordism theory.'' \emph{Bulletin of the American
 Mathematical Society}, 75(6), 1293--1298.
\item
 Ranicki, A. (Ed.). (1999). \emph{The Novikov Conjectures on Higher
 Signatures and K-Theory}. London Mathematical Society Lecture Note
 Series, Cambridge University Press.
\item
 Veselov, A. P., \& Novikov, S. P. (1984). ``Finite-zone
 two-dimensional Schrödinger operators. Potential operators.''
 \emph{Doklady Akademii Nauk SSSR}, 279(1), 20--24.
\end{enumerate}

\begin{center}\rule{0.5\linewidth}{0.5pt}\end{center}

\emph{(End of report)}
